%---------------------------------------------------------------
%   CHAPTER 1 
%---------------------------------------------------------------
\chapter{Interfejs mózg-komputer}
\label{rozdzial1}

Od najdawniejszych czasów człowiek starał się poszerzać swoje możliwości poznawcze i praktyczne za pomocą narzędzi. Pierwsze z nich - jak liczydła czy kamienne tablice - służyły do wspomagania podstawowych operacji arytmetycznych. Z biegiem stuleci te proste urządzenia przekształciły się w coraz bardziej złożone maszyny obliczeniowe, aż po koncepcje maszyn programowalnych. U podstaw tego rozwoju leżała nie tylko potrzeba zwiększenia mocy obliczeniowej, lecz także dążenie do uproszczenia komunikacji między człowiekiem a narzędziem.

Za prekursora tej idei uważa się Charlesa Babbage'a, który w XIX wieku zaprojektował maszynę analityczną - pierwszą koncepcję programowalnego komputera \cite{campbellkelly2009}. W jego projekcie dane oraz instrukcje miały być wprowadzane za pomocą kart perforowanych, a wyniki - drukowane lub zapisywane na kartach perforowanych w celu możliwości dalszego przetwarzania. Choć maszyna Babbage'a nigdy nie została w pełni zrealizowana, wyznaczyła kierunek dla późniejszego rozwoju informatyki: automatyzacji przetwarzania informacji oraz poszukiwania sposobów na bardziej bezpośrednią interakcję z maszyną. Pierwsze komputery XX wieku były urządzeniami zarezerwowanymi dla specjalistów. Programowanie wymagało manualnego łączenia przewodów i ustawiania przełączników. Podstawową barierą komunikacyjną stanowiła różnica między symbolicznym językiem człowieka a binarną reprezentacją danych w komputerze. Każda interakcja była czasochłonna i obca ludzkiej intuicji. W drugiej połowie lat 50. XX wieku pojawiły się pierwsze próby połączenia komputerów z maszynami do pisania. Jednym z pionierskich rozwiązań był IBM 610, który jako jeden z pierwszych komputerów wykorzystywał klawiaturę do wprowadzania danych \cite{campbellkelly2009}. W połączeniu z interfejsem tekstowym, doprowadziło to do popularyzacji klawiatur jako urządzeń wejściowych oraz rozwoju języków programowania.

W miarę jak komputery stawały się coraz bardziej dostępne, wzrastało zapotrzebowanie na bardziej intuicyjne sposoby interakcji z nimi. Klawiatura i interfejs tekstowy umożliwiły użytkownikowi komunikację za pomocą języka pisanego, lecz wciąż była to forma symboliczna, daleka od naturalnych działań człowieka w świecie fizycznym. Przełom nastąpił w latach 60., kiedy amerykański badacz Douglas Engelbart opracował w Stanford Research Institute pierwszy prototyp myszy komputerowej \cite{engelbart1968}. Urządzenie wykonane z drewna, wyposażone w dwa metalowe kółka, pozwalało użytkownikowi wskazywać i manipulować obiektami na ekranie w sposób przypominający ruch ręki po powierzchni stołu. W 1968 roku Engelbart zaprezentował swój wynalazek podczas słynnej prezentacji „The Mother of All Demos”, w której po raz pierwszy pokazano także działające okna, hiperłącza i współpracę online w czasie rzeczywistym. Demonstracja ta stała się punktem zwrotnym w historii interfejsów użytkownika, wyznaczając kierunek rozwoju ku środowiskom graficznym. Idea Engelbarta, rozwinięta później przez ośrodki badawcze Xerox PARC \cite{myers1998} i komercyjnie wdrożona przez firmę Apple w latach 80., doprowadziła do powstania współczesnych interfejsów graficznych (GUI). Od tamtego momentu zamiast tylko wydawać abstrakcyjne polecenia w terminalu człowiek mógł również manipulować obiektami na ekranie.

% 1.1
\section{Komunikacja z komputerem w dzisiejszych czasach}

Współczesny świat zdominowany jest przez urządzenia komputerowe, które stały się nieodłącznym elementem codziennego życia. Przełomowym momentem w rozwoju technologii interakcji człowieka z maszyną był rok 2007, kiedy firma Apple zaprezentowała iPhone'a - pierwszy smartfon w pełni oparty na ekranie dotykowym \cite{isaacson2011}. Choć sama idea dotykowej obsługi urządzeń nie była nowa, to właśnie w tym okresie rozpoczęła się jej masowa popularyzacja. Już wcześniej ekrany dotykowe znajdowały zastosowanie w terminalach, bankomatach i specjalistycznych panelach sterujących, jednak dopiero rozwój elektroniki użytkowej uczynił je standardem. Warto dodać, że początki technologii ekranów dotykowych sięgają lat sześćdziesiątych XX wieku \cite{johnson1965}.

W konstrukcji ekranów dotykowych stosuje się kilka odmiennych zasad działania. Najstarsze z nich wykorzystywały zjawisko zmiany oporu elektrycznego pomiędzy przezroczystymi warstwami przewodzącymi. Rozwiązania tego typu, nazywane ekranami rezystancyjnymi, pozwalały na obsługę przy użyciu dowolnego przedmiotu, w tym rysika, lecz cechowały się mniejszą precyzją i niższą odpornością mechaniczną. W późniejszych latach opracowano technologie oparte na przerwaniu wiązki podczerwieni emitowanej przez diody rozmieszczone na krawędziach ekranu oraz na zaburzeniu fali akustycznej rozchodzącej się po jego powierzchni (technologia \textit{Surface Acoustic Wave}, SAW). Technologia detekcji w podczerwieni znalazła zastosowanie nie tylko w ekranach dotykowych, ale również w innych interfejsach użytkownika, takich jak kontroler ruchu konsoli Nintendo Wii. W tym rozwiązaniu optyczny czujnik w pilocie odczytywał pozycję dwóch diod podczerwonych w tzw. sensor barze, a dane z czujników inercyjnych (akcelerometrów i żyroskopu) umożliwiały precyzyjne odwzorowanie ruchu w trzech wymiarach. Największą popularność zdobyły jednak ekrany pojemnościowe, które działają na zasadzie zmiany pola elektrostatycznego w miejscu kontaktu z przewodzącym materiałem, takim jak ludzka skóra. Umożliwiają one precyzyjną i płynną obsługę, a także wykrywanie wielu punktów dotyku jednocześnie (funkcja \textit{multi-touch}). Ich podstawową wadą jest brak reakcji na przedmioty nieprzewodzące prądu, co uniemożliwia obsługę w rękawiczkach lub za pomocą tradycyjnego rysika. Pomimo wyższych kosztów produkcji, technologia pojemnościowa niemal całkowicie zastąpiła wcześniejsze rozwiązania i obecnie stanowi standard w smartfonach, tabletach oraz panelach sterujących.

Kolejnym etapem rozwoju metod interakcji człowieka z komputerem stały się interfejsy bezdotykowe, wykorzystujące rozpoznawanie gestów, mowy, ruchu ciała lub kierunku spojrzenia użytkownika. Rozwiązania tego typu umożliwiają komunikację z urządzeniem bez fizycznego kontaktu, co znacząco zwiększa intuicyjność obsługi i poszerza zakres możliwych zastosowań. Przykładem takiej technologii jest system \textit{Microsoft Kinect}, który po raz pierwszy został zaprezentowany w 2010 roku jako urządzenie peryferyjne dla konsoli Xbox 360. System ten wykorzystywał zestaw kamer, czujników głębi oraz mikrofonów do analizy pozycji ciała, gestów i poleceń głosowych użytkownika w czasie rzeczywistym. Dzięki temu możliwe było sterowanie grami i aplikacjami za pomocą ruchu rąk, tułowia czy komend głosowych, bez konieczności używania tradycyjnego kontrolera. Wprowadzenie technologii głębi 3D umożliwiło również śledzenie przestrzenne wielu użytkowników jednocześnie, co stanowiło istotny krok w kierunku bardziej naturalnej interakcji z systemami komputerowymi.

Najbardziej naturalną formą komunikacji człowieka z otoczeniem jest mowa, dlatego też próby jej wykorzystania jako środka sterowania komputerem rozpoczęto już w latach pięćdziesiątych XX wieku. W 1952 roku K. H. Davis, R. Biddulph i S. Balashek opracowali pierwszy system zdolny do rozpoznawania wypowiadanych cyfr \cite{davis1952}. W kolejnych dekadach rozwój tej technologii koncentrował się na systemach ograniczonych słownikowo lub dostosowanych do konkretnego użytkownika, co pozwalało zwiększyć dokładność rozpoznawania w warunkach ograniczonej mocy obliczeniowej. Rozpoznawanie mowy jest zagadnieniem złożonym, ponieważ wymaga analizy sygnału akustycznego oraz jego interpretacji językowej. Zmienność głosu, akcentu, tempa mówienia czy obecność zakłóceń sprawiają, że stworzenie uniwersalnego systemu zdolnego do rozumienia mowy potocznej stanowi duże wyzwanie. Współczesne rozwiązania wykorzystują jednak zaawansowane modele akustyczne i językowe oparte na sztucznych sieciach neuronowych, co pozwoliło znacząco zwiększyć dokładność i naturalność interakcji głosowej.

Technologie rozpoznawania mowy znajdują zastosowanie przede wszystkim w sterowaniu urządzeniami oraz w transkrypcji tekstu mówionego. Coraz częściej pełnią także rolę narzędzia komunikacji, umożliwiając porozumiewanie się osób mówiących w różnych językach w czasie rzeczywistym. Powszechne stały się systemy asystentów głosowych, takich jak Siri, Alexa czy Google Assistant, które pozwalają wykonywać polecenia, wyszukiwać informacje i sterować inteligentnymi urządzeniami domowymi przy użyciu komend słownych. Rozwiązania te stanowią ważny krok w kierunku bardziej naturalnej komunikacji z komputerem, przybliżając ją do sposobu, w jaki ludzie porozumiewają się ze sobą.

W idealnym świecie człowiek nie musiałby wydawać poleceń komputerowi - maszyna sama wiedziałaby, co chcemy osiągnąć, jakby czytała nasze myśli. Choć wizja ta pozostaje w sferze science fiction, jej częściowe realizacje istnieją już w świecie rzeczywistym. Obecnie największym ograniczeniem interfejsów BCI jest stosunkowo niska szybkość komunikacji. Mimo to systemy te znalazły zastosowanie w psychologii, technikach multimedialnych, wojskowości, a przede wszystkim w medycynie. Tam stanowią często jedyny sposób komunikacji ze światem zewnętrznym dla osób dotkniętych poważnymi chorobami neurologicznymi - takimi jak stwardnienie zanikowe boczne, udary mózgu czy urazy rdzenia kręgowego - które stopniowo ograniczają kontrolę nad ruchem ciała. W takich przypadkach interfejs mózg-komputer pozwala nie tylko na podstawową komunikację, ale również na sterowanie protezami, wózkiem inwalidzkim czy inteligentnym domem.

% 1.2
\section{Zastosowania interfejsów mózg-komputer}

Interfejsy mózg-komputer znajdują zastosowanie w wielu dziedzinach. Ich rozwój koncentruje się głównie na umożliwieniu osobom z niepełnosprawnościami komunikacji oraz kontroli urządzeń. W medycynie systemy BCI są wykorzystywane do przywracania funkcji ruchowych u pacjentów z paraliżem, umożliwiając sterowanie protezami kończyn lub wózkami inwalidzkimi za pomocą myśli \cite{hochberg2012}. Ponadto BCI jest wykorzystywane w celu wspomagania komunikacji osób z zespołem zamknięcia (locked-in syndrome), które nie mogą poruszać się ani mówić. Dzięki analizie sygnałów EEG możliwe jest rozpoznawanie intencji ruchu, co pozwala na wykonywanie prostych czynności, takich jak poruszanie kursorem lub wybór znaków na wirtualnej klawiaturze. 

Jednym z wielu obszarów zastosowań jest neurorehabilitacja, gdzie BCI wspomagają procesy odzyskiwania funkcji motorycznych po udarach mózgu lub urazach rdzenia kręgowego. Systemy te wykorzystują neuroplastyczność mózgu, umożliwiając pacjentom ćwiczenie kontrolowanych ruchów poprzez wizualizację i sprzężenie zwrotne \cite{dobkin2007}. W praktyce BCI często integruje się z egzoszkieletami lub systemami funkcjonalnej stymulacji elektrycznej (FES), co pozwala na bezpośrednie wspomaganie wykonywanych ruchów.
Badania kliniczne potwierdzają skuteczność takiego podejścia - w badaniu Frolova i in. zastosowanie egzoszkieletu sterowanego poprzez wyobrażenie ruchu doprowadziło do istotnej poprawy funkcji ręki u pacjentów po udarze w porównaniu z grupą kontrolną \cite{frolov2017}. Podobne wyniki uzyskano w eksperymencie Bundy'ego i współpracowników, gdzie uczestnicy sterowali egzoszkieletem kończyny górnej za pomocą sygnałów EEG z nieuszkodzonej półkuli, co przełożyło się na wzrost aktywności korowej i poprawę wyników testów motorycznych \cite{bundy2017}. Takie zamknięte pętle treningowe mogą zwiększać aktywność korową i przyspieszać odzyskiwanie kontroli motorycznej, szczególnie przy intensywnym i spersonalizowanym treningu. Efektywność terapii zależy jednak od stopnia uszkodzenia, czasu od urazu oraz dopasowania protokołu rehabilitacyjnego.

Coraz większe znaczenie zyskują także zastosowania poznawcze i badawcze, obejmujące analizę uwagi, zmęczenia czy emocji użytkownika. W takich przypadkach BCI służy do oceny stanu psychofizjologicznego w czasie rzeczywistym, co może być użyteczne np. w transporcie, wojsku lub podczas złożonych zadań wymagających wysokiej koncentracji.

% 1.3
\section{Struktura interfejsu mózg-komputer}

Podstawowa struktura interfejsu mózg-komputer (BCI) została opracowana w latach 70. XX wieku wraz z pracami Jacquesa Vidala, który wprowadził pojęcie BCI i przeprowadził pierwsze eksperymenty umożliwiające sterowanie kursorem lub robotem przy użyciu sygnałów EEG \cite{vidal1973,vidal1977}. System BCI składa się z kilku kluczowych etapów przetwarzania sygnału. Pierwszym jest rejestracja aktywności mózgu, najczęściej przy użyciu EEG. Zarejestrowane sygnały wymagają wstępnego oczyszczenia - filtracji zakłóceń fizjologicznych i technicznych - oraz ekstrakcji cech, czyli wyodrębnienia parametrów najbardziej istotnych do interpretacji intencji użytkownika. W tym celu stosuje się filtry dolno- i górno- przepustowe oraz filtry pasmozaporowe (notch filter). Kolejnym krokiem jest klasyfikacja tych cech przy użyciu algorytmów uczenia maszynowego, który pozwala na przypisanie im określonych komend lub decyzji. Ostatni moduł systemu tłumaczy sygnały mózgowe na konkretne działania, takie jak ruch kursora, sterowanie protezą czy wózkiem inwalidzkim. Na rysunku \ref{fig:BCI_structure} przedstawiono schemat działania interfejsu mózg-komputer. Centralnym elementem całego systemu pozostaje komputer, który analizuje, klasyfikuje sygnały oraz generuje sprzężenie zwrotne, tworząc zamkniętą pętlę komunikacyjną między mózgiem a urządzeniem wykonawczym \cite{wolpaw2002,wolpaw2012}.

\begin{figure}[!ht]
    \centering
    \includegraphics[width=0.8\textwidth]{diagram_bci.png}
    \caption{Schemat struktury interfejsu mózg-komputer (Źródło: opracowanie własne na podstawie \cite{kolodziej2011przetwarzanie})}
    \label{fig:BCI_structure}
\end{figure}

% 1.4
\section{Metody badania aktywności mózgu}

Aby interfejs mózg-komputer mógł działać, niezbędne jest poznanie i rejestracja aktywności elektrycznej mózgu. Istnieje wiele technik badania pracy mózgu, różniących się stopniem inwazyjności, dokładnością i zastosowaniami \cite{nicolas-alonso2012,he2019}. Metody inwazyjne obejmują implantowane elektrody domózgowe (IntraCortical Recordings, IR), które pozwalają na bezpośrednie odczytywanie sygnałów neuronów z bardzo wysoką rozdzielczością czasową i przestrzenną. Ich zastosowanie ogranicza jednak wysoki stopień ingerencji chirurgicznej i związane z tym ryzyko. Metody częściowo inwazyjne, takie jak elektrokortykografia (Electrocorticography, ECoG), wykorzystują elektrody umieszczane na powierzchni kory mózgowej. Dają one lepszą jakość sygnału niż metody nieinwazyjne, przy mniejszym ryzyku niż pełne implanty. Metody nieinwazyjne nie wymagają ingerencji chirurgicznej i są bezpieczne do wielokrotnych pomiarów u zdrowych osób. Należą do nich:

\begin{itemize}
    \item Magnetoencefalografia (MEG) - rejestruje pola magnetyczne generowane przez aktywność neuronów
    \item Funkcjonalny rezonans magnetyczny (fMRI) - mierzy zmiany przepływu krwi w mózgu jako wskaźnik aktywności neuronów
    \item Tomografia optyczna (NIRS) - wykorzystuje bliską podczerwień do pomiaru przepływu krwi i saturacji tlenu w tkankach mózgowych.
    \item Pozytronowa tomografia emisyjna (PET) - pozwala śledzić metabolizm i przepływ krwi
    \item Elektroencefalografia (EEG) - rejestruje elektryczną aktywność mózgu za pomocą elektrod umieszczonych na skórze głowy
\end{itemize}

Choć istnieje kilka metod nieinwazyjnego badania aktywności mózgu, tylko niektóre nadają się do zastosowania w systemach BCI codziennego użytku. Techniki takie jak MEG, PET, fMRI czy tomografia optyczna oferują precyzyjną lokalizację aktywnych obszarów mózgu, jednak są kosztowne, wymagają skomplikowanego sprzętu i specjalistycznej obsługi. Co więcej, mierzą pośrednie efekty aktywności neuronów, takie jak zmiany przepływu krwi, które zachodzą znacznie wolniej niż rzeczywiste impulsy nerwowe, ograniczając szybkość interakcji. Elektroencefalografia za to, przeciwnie, pozwala rejestrować aktywność mózgu w czasie rzeczywistym, jest stosunkowo tania i łatwa w użyciu. Dzięki temu systemy BCI oparte na EEG mogą działać w różnych warunkach, a po krótkim przeszkoleniu użytkownika - nawet w domu. To czyni EEG najbardziej praktycznym wyborem do interfejsów mózg-komputer. Wadą tego rozwiązania jest niski stosunek sygnału do szumu, wynikający z podatności na zakłócenia fizjologiczne (np. napięcie mięśniowe, ruchy oczu) i techniczne (zakłócenia sieci elektrycznej).
Sygnał EEG odzwierciedla zbiorową aktywność dużych grup neuronów, co sprawia, że pomiar jest obarczony naturalnymi ograniczeniami dokładności. W odróżnieniu od pomiarów technicznych, gdzie obserwowany parametr istnieje wprost, w biologii sam obiekt pomiaru bywa mniej jednoznaczny. Z tego powodu pojedyncza rejestracja EEG rzadko dostarcza pełnej informacji, a wiarygodne wyniki wymagają wielokrotnych powtórzeń i zastosowania analizy statystycznej. Pomimo tych ograniczeń, sygnał EEG dostarcza cennych danych o stanie funkcjonalnym mózgu. Analiza jego struktury czasowo-częstotliwościowej pozwala wyróżnić charakterystyczne wzorce aktywności neuronalnej, zwane rytmami EEG, które stanowią podstawę interpretacji i wykorzystania sygnału w systemach BCI.

% 1.5
\section{Pomiar sygnału EEG}
Do rejestracji sygnału EEG stosuje się zestaw elektrod przytwierdzanych do skóry głowy, zazwyczaj z użyciem przewodzącego żelu. Elektrody te wychwytują aktywność elektryczną mózgu, która następnie jest wzmacniana i przetwarzana w formę cyfrową przez elektroencefalograf. Urządzenie to mierzy różnice potencjałów między elektrodami umieszczonymi na głowie. Sygnały te są bardzo słabe - ich amplituda nie przekracza kilkudziesięciu mikrovoltów - dlatego niezbędne jest zastosowanie wzmacniaczy. Po wzmocnieniu sygnał jest próbkowany i zapisywany w postaci cyfrowej, co umożliwia jego dalszą analizę.

Rozmieszczenie elektrod na głowie odbywa się według tzw. międzynarodowego systemu 10-20, opracowanego przez Herberta Jaspera w 1958 roku. System ten określa pozycje elektrod w oparciu o odległości pomiędzy charakterystycznymi punktami anatomicznymi czaszki: nasadą nosa (\textit{nasion}), guzowatością potyliczną (\textit{inion}) oraz punktami przeduszno-skroniowymi \cite{jasper1958}. Odległości między sąsiednimi elektrodami stanowią odpowiednio 10$\%$ lub 20$\%$ całkowitej długości wybranego wymiaru głowy, co zapewnia powtarzalność i porównywalność pomiarów między badaniami. Każda elektroda oznaczona jest literą odpowiadającą regionowi mózgu, nad którym się znajduje (F - czołowy, C - centralny, P - ciemieniowy, O - potyliczny, T - skroniowy) oraz cyfrą: nieparzyste numery odnoszą się do lewej półkuli, a parzyste do prawej. Punkt środkowy oznacza się literą „z” (od ang. \textit{zero line}), np. Fz lub Cz. W praktyce standardowy zestaw EEG obejmuje od 19 do 21 elektrod, jednak w badaniach naukowych często stosuje się systemy o większej gęstości (np. 10-10 lub 10-5), zawierające nawet kilkaset elektrod, co pozwala na dokładniejsze odwzorowanie przestrzennego rozkładu aktywności mózgu. Na rysunku \ref{fig:elektrody_na_glowie} przedstawiono schemat rozmieszczenia elektrod w standardowym systemie 10-20.

\begin{figure}[!ht]
    \centering
    \includegraphics[width=0.4\textwidth]{system10-20.png}
    \caption{Rozmieszczenie elektrod na powierzchni głowy w standardowym systemie 10-20 stosowanym w elektroencefalografii}
    \label{fig:elektrody_na_glowie}
\end{figure}

Elektrody EEG mogą być wykonane z różnych materiałów przewodzących, takich jak srebro, chlorek srebra (Ag/AgCl) czy złoto. Najczęściej stosuje się elektrody Ag/AgCl ze względu na ich stabilność potencjału i odporność na dryft napięcia. W zależności od sposobu montażu wyróżnia się elektrody przyklejane, żelowe, suche oraz półsuche. Elektrody suche eliminują konieczność użycia żelu przewodzącego, co zwiększa wygodę użytkownika, ale często wiąże się z większym szumem kontaktowym i mniejszą jakością rejestracji.

Typowy system EEG obejmuje również elektrodę odniesienia (\textit{reference electrode}) oraz elektrodę uziemiającą (\textit{ground electrode}). Elektroda odniesienia stanowi punkt odniesienia potencjałów mierzonych przez pozostałe elektrody i może być umieszczona np. na małżowinie usznej lub w punkcie Cz. Sposób wyboru elektrody odniesienia ma istotny wpływ na kształt i interpretację zarejestrowanego sygnału, dlatego w analizie EEG często stosuje się techniki przeliczeniowe, np. odniesienie wspólne lub tzw. \textit{average reference}, w którym od każdego kanału odejmuje się średnią wartość sygnałów ze wszystkich elektrod. Współczesne elektroencefalografy umożliwiają rejestrację wielu kanałów jednocześnie (od 8 do nawet 256), a także integrację z innymi systemami pomiarowymi, np. śledzenia wzroku, EMG czy EKG. Urządzenia te coraz częściej wykorzystują technologie bezprzewodowe, co zwiększa mobilność i umożliwia prowadzenie pomiarów poza laboratorium. Dzięki postępowi w miniaturyzacji elektroniki, przenośne zestawy EEG stają się coraz powszechniejsze również w zastosowaniach domowych, edukacyjnych i rozrywkowych.

% 1.6
\section{Rytmy EEG}

Rejestrowany sygnał EEG ma bardzo małą amplitudę - zwykle od kilku do kilkudziesięciu $\mu$V i jest silnie narażony na zakłócenia. Źródłem szumów mogą być zarówno czynniki zewnętrzne, jak zakłócenia sieciowe o częstotliwości 50 Hz (lub 60Hz) czy urządzenia elektryczne w otoczeniu, jak i wewnętrzne, takie jak ruchy mięśni twarzy, napięcie mięśni czoła, ruchy gałek ocznych czy sygnały z serca. W praktyce zakłócenia te często osiągają amplitudy wielokrotnie większe niż właściwy sygnał mózgowy, co wymaga stosowania zaawansowanych metod filtracji i przetwarzania. Mimo to, analiza EEG pozwala na identyfikację charakterystycznych wzorców aktywności mózgu, zwanych rytmami lub falami. Są to sygnały o względnie regularnej strukturze i określonym zakresie częstotliwości, które odzwierciedlają stan funkcjonalny mózgu. Wyróżnia się pięć głównych rytmów: delta, theta, alfa, beta i gamma. Ich parametry mogą się różnić w zależności od wieku, stanu psychicznego oraz fizjologicznego badanego. Na rysunku \ref{fig:syngaly_eeg} przedstawiono przykładowe sygnały poszczególnych rytmów EEG.

\paragraph{Rytm delta (0,5-4 Hz)}

Rytmy delta występują w najniższym zakresie częstotliwości. Są charakterystyczne dla snu głębokiego i stanów utraty świadomości. U dorosłych pojawienie się fal delta w stanie czuwania może wskazywać na nieprawidłowości w funkcjonowaniu ośrodkowego układu nerwowego. W zapisie EEG dzieci rytmy delta są jednak zjawiskiem fizjologicznym. Typowa amplituda fal delta wynosi około 50 µV, a ich obecność obserwuje się zwykle na obszarze czołowo-centralnym i w stanie relaksu lub snu.

\paragraph{Rytm theta (4-7 Hz)}

Rytmy theta wiążą się z obniżonym poziomem czuwania, sennością, marzeniami sennymi oraz wczesnymi fazami snu. Występują także podczas stanów głębokiego relaksu, medytacji czy chwilowego oderwania uwagi od bodźców zewnętrznych. U dorosłych duży udział fal theta w stanie aktywnego czuwania może wskazywać na zaburzenia neurologiczne lub nadmierne zmęczenie. Amplituda tego rytmu nie przekracza zazwyczaj 30 µV. W zapisie EEG fale theta pojawiają się głównie w odprowadzeniach środkowych.

\paragraph{Rytm alfa (8-13 Hz)}

Rytmy alfa zostały opisane po raz pierwszy przez Hansa Bergera w 1930 roku i są uznawane za podstawowy wzorzec aktywności mózgu u dorosłych. Występują głównie w stanie spoczynku, przy zamkniętych oczach i braku intensywnego skupienia. Zanik rytmu alfa następuje po otwarciu oczu lub w sytuacji silnego pobudzenia psychicznego. Sygnał ten jest najlepiej widoczny w okolicach potylicznych, czyli nad korą wzrokową. Częstotliwość rytmu alfa mieści się w przedziale 8-13 Hz, a amplituda wynosi zwykle 20-100 µV. Całkowity brak tego rytmu u osoby dorosłej w stanie relaksu może wskazywać na zaburzenia neurologiczne.

\paragraph{Rytm beta (13-30 Hz)}

Rytmy beta obejmują fale o wyższych częstotliwościach i mniejszej amplitudzie, które dominują podczas aktywności umysłowej, koncentracji, analizy informacji oraz w trakcie ruchu. Występują głównie w przednich częściach mózgu, szczególnie w rejonach czołowo-centralnych. Rytm beta wykazuje znaczną zmienność amplitudy, co ogranicza jego zastosowanie diagnostyczne. Obserwuje się także asymetrię międzypółkulową, która może mieć znaczenie kliniczne - np. słabsza obecność rytmu beta w jednej półkuli bywa powiązana z zaburzeniami funkcjonalnymi. Część leków, zwłaszcza barbiturany, zwiększa amplitudę i udział fal beta w zapisie EEG.

\paragraph{Rytm gamma (30-70 Hz)}

Fale gamma występują w najwyższym zakresie częstotliwości i wiązane są z procesami poznawczymi wyższego rzędu: percepcją, zapamiętywaniem, integracją informacji sensorycznych oraz aktywnością ruchową. Aktywność gamma odzwierciedla szybkie procesy synchronizacji między różnymi obszarami mózgu i jest uznawana za wskaźnik intensywnego przetwarzania informacji. Choć fale te mają niewielką amplitudę, ich analiza dostarcza cennych danych o mechanizmach świadomości i uwagi.

\begin{figure}[!ht]
    \centering
    \includegraphics[width=0.8\textwidth]{rytmy_eeg.png}
    \caption{Próbki sygnałów poszczególnych rytmów EEG \cite{wikipedia_fale_mozgowe}}
    \label{fig:syngaly_eeg}
\end{figure}